% Source: https://artofproblemsolving.com/community/c3h1931055_binary_numbers_and_polygons
% Stable identity: aops-topic-1931055-q02
There are $64-24=40$ hexagons. Count the polygon edges: $24\cdot5+40\cdot6=120+240=360$. Every seam joins exactly two polygon edges, one from each of two neighbouring polygons, so each seam is counted twice:

\begin{asy}
usepackage("amsmath");
size(4cm);
defaultpen(fontsize(10pt));
pair[] H={(0,0),(1,0),(1.5,sqrt(3)/2),(1,sqrt(3)),(0,sqrt(3)),(-0.5,sqrt(3)/2)};
pair[] P={(0,0),(1,0)};
for(int k=2;k<5;++k) P.push(P[k-1]+dir(-72*(k-1)));
draw(H[0]--H[1]--H[2]--H[3]--H[4]--H[5]--cycle,linewidth(1));
draw(P[0]--P[1]--P[2]--P[3]--P[4]--cycle,linewidth(1));
draw(H[0]--H[1],red+linewidth(2));
label("hexagon",(0.5,sqrt(3)/2));
label("pentagon",(0.5,-0.69));
label("one seam",(1.25,0),E,red);
label("\mbox{}",(min(currentpicture,user=true)+max(currentpicture,user=true))/2,black); // keeps the text after the figure black when LaTeX typesets the labels inline
\end{asy}

\[
\text{number of seams}=\dfrac{360}{2}=\boxed{180}.
\]
(Remark: a ball sewn from pentagons and hexagons with exactly three polygons at every corner must have exactly $12$ pentagons, by Euler's formula; the ball described here must have some corners where more polygons meet. This does not affect the count of seams.)
